\subsection{Artefatti per la simulazione dei regimi operativi}
\label{sec:artefatti-regimi-operativi}

Gli artefatti della simulazione dei regimi operativi sono separati dal
runtime router. Essi consumano CSV e JSON gia prodotti dal benchmark
R-D-E e generano tabelle, dati di plotting e figure derivate per la
tesi. Nessun file Python deve essere copiato nel progetto LaTeX: nella
tesi entrano soltanto figure PNG, tabelle LaTeX e testo descrittivo.

\begin{table}[t]
\centering
\caption{Artefatti per la simulazione dei regimi operativi e la loro integrazione nella tesi.}
\label{tab:operational-regime-thesis-artifacts}
\begin{tabular}{p{0.30\linewidth} p{0.58\linewidth}}
\toprule
Artefatto & Ruolo nella pipeline di tesi \\
\midrule
\texttt{operational\_regime\_simulation.py} &
Genera gli artefatti scientifici primari della simulazione, cioe CSV e
JSON con decisioni, summary, switch analysis, quality contract e safety
gate. Non genera PNG nel path principale. \\
\texttt{operational\_regime\_plots.py} &
Legge i CSV finali della simulazione e produce figure PNG derivate. Il
plotter e rieseguibile indipendentemente e un errore di rendering non
invalida i CSV/JSON scientifici. \\
\texttt{operational\_regime\_diagnostics.py} &
Diagnostica la consistenza tra metrica di qualita, oracle, objective,
regret e no-leakage. Il suo ruolo e interpretativo e read-only. \\
\texttt{operational\_regime\_switch\_summary.csv} &
Riassume, per regime e quality floor, quando vince il miglior candidato
classico o neurale e con quale switch reason. Alimenta la Tabella
\ref{tab:ssimulacra2-switch-summary}. \\
\texttt{switch\_reason\_by\_rate\_weight.csv} &
Fornisce i dati per visualizzare come cambiano le switch reason al
crescere del peso sul rate. Alimenta la Figura
\ref{fig:switch-reason-rate-weight}. \\
\texttt{quality\_gate\_comparison\_summary.csv} &
Confronta la policy predittiva senza gate e la variante shadow con
expected-quality gate. Alimenta la Tabella
\ref{tab:expected-quality-gate-summary}. \\
\bottomrule
\end{tabular}
\end{table}

Questa organizzazione distingue gli artefatti scientifici primari dai
rendering derivati. I CSV/JSON sono la base riproducibile dei risultati;
le figure PNG sono viste grafiche rigenerabili a partire dagli stessi
dati.
